\section{本章小结}
\label{sec:odesens-summary}

章首的问题是：把$x_0$改成$0.26$，或把增长率改$20\%$，$t=\ln3$时的值会差多少？现在有三种回答。Gronwall上界告诉我们，两种改变的影响分别不超过$0.01\,\mathrm{e}^{\ln3}=0.03$和$0.1$。变分方程给出一阶值$\tfrac43\times0.01=0.0133$与$0.0549$，后者与真实差$0.0547$很接近。若只需要一个目标对很多参数的梯度，伴随方程用一次反向求解就能得到。

四个工具输出的量不同，见表~\ref{tab:odesens-tools}。

\begin{BookInlineTable}
  \centering
  \caption{本章工具的对比。结论都针对给定解附近的紧时间段。}
  \label{tab:odesens-tools}
  \begin{tabular}{p{16mm}p{24mm}p{27mm}p{28mm}}
    \toprule
    工具 & 输出 & 主要条件 & 不能回答 \\
    \midrule
    Gronwall估计 & 曲线距离的上界 & 状态Lipschitz，系数可积 & 差的方向；长期趋势 \\
    变分方程 & 一阶响应$\bm{S}$、$\bm{P}$ & 对状态、参数$C^1$，导数有可积上界 & 大扰动；非光滑处 \\
    流映射 & 复合、逆、体积变化 & 解唯一，$C^1$ & 定义域外的解；长期轨道 \\
    伴随方程 & 目标函数的梯度 & 同变分方程，目标$C^1$ & 整个$\bm{P}(t)$；控制函数的最优性 \\
    \bottomrule
  \end{tabular}
\end{BookInlineTable}

几个相近的概念要分清。连续依赖只说解差随数据差趋于零；Lipschitz依赖给出速度，需要右端局部Lipschitz；可微依赖给出一阶系数，需要右端$C^1$，这三者一层比一层强，且不能由弱推强，$\dot{x}=-x\ln\lvert x\rvert$和$\dot{x}=\lvert x\rvert$分别是分界的反例。积分范数适合度量右端关于输入Lipschitz时的输入差，非线性增长时要限制输入的取值或改用更强的范数。残差和放大因子分别来自求解器和方程。

所有结论都是有限时间的。常数$\mathrm{e}^{\Lambda(T)}$可以很保守，也不说明长时间下解是靠拢还是分离，长期稳定需要卷九的额外结构和条件。第~\ref{chap:ode-numerical-methods}章的全局误差分析将回用本章的扰动估计，第~\ref{chap:transport-equations}章的输运方程回用流映射和体积变化。
